% related-work.tex — HapGraph Related Work
% Target: 500+ words, 10+ citations
% Must POSITION against prior work, not just list it

\subsection*{Admixture Graph Inference}

The problem of inferring admixture graphs from population genetic data has
a rich history rooted in the development of F-statistics.
\citet{Reich2009} introduced the $F_3$ and $F_4$ statistics as formal tests
for tree-inconsistency, providing the first statistical framework for
detecting admixture in a graph context.
\citet{Patterson2012} formalised the full suite of F-statistics and their
bias-corrected estimators, proving that negative $F_3(C;\,A,B)$ is a
statistically unambiguous signal of admixture when $C$ lies on an internal
branch between $A$ and $B$; this foundation underlies both \hapgraph{} and
all subsequent graph-inference tools.

\treemix{} \citep{Pickrell2012} was the first fully automated admixture
graph inference tool, placing the problem in a maximum-likelihood framework
and using a greedy search over admixture edge proposals.
\hapgraph{}'s Stage I topology search is conceptually similar to \treemix{},
using the same F-statistic covariance structure, but differs in two key
respects: the stopping criterion is BIC-based (scale-adaptive) rather than
a fixed residual threshold, and the pipeline subsequently estimates
admixture timing — a capability entirely absent from \treemix{}.

\textsc{qpGraph} \citep{Haak2015} and its successor \textsc{ADMIXTOOLS2}
\citep{Maier2023} extended F-statistic graph inference to support user-specified
topologies with formal $F$-ratio optimisation, enabling precise proportion
estimation under a fixed graph.
\textsc{ADMIXTOOLS2} also implements an automated search over graph space
\citep{Maier2023}.
These tools do not model admixture timing, and \hapgraph{}'s topology-independent
$F_3$ method-of-moments estimator is conceptually related to \textsc{qpGraph}'s
proportion fitting but does not require a correctly specified topology.

\textsc{AdmixtureBayes} \citep{Nielsen2023} represents the current state of
the art in Bayesian admixture graph inference, using a reversible-jump MCMC
sampler to jointly explore graph topology and proportion space with full
posterior uncertainty quantification.
Its treatment of topology uncertainty is more principled than \hapgraph{}'s
two-stage approach.
However, \textsc{AdmixtureBayes} does not model admixture timing, and its
RJMCMC sampler becomes computationally intractable for large numbers of
populations, making it unsuitable for datasets with $N \gg 10$ populations.
\hapgraph{} trades formal joint topology uncertainty for scalability and a
timing estimation capability that \textsc{AdmixtureBayes} lacks.

\subsection*{Admixture Timing}

\alder{} \citep{Loh2013} estimates admixture timing from the decay of
weighted linkage disequilibrium (LD) between ancestry-informative markers,
exploiting the fact that admixture-induced LD decays at a rate of
$\sim\!1/T$ per generation.
\alder{} is the standard timing complement to \treemix{}; the two tools
are typically applied sequentially, with \treemix{} providing the topology
and \alder{} providing timing.
\hapgraph{} replaces this two-step workflow with a single integrated analysis
and uses IBD segment lengths rather than LD decay as the timing signal, which
is complementary: IBD is more sensitive to recent admixture ($T < 50$
generations) while LD-based methods can in principle reach deeper events.

\textsc{GLOBETROTTER} \citep{Hellenthal2014} jointly infers admixture sources
and timing from chromosome painting coancestry curves, making it conceptually
the closest existing tool to \hapgraph{} in its goal of joint source and timing
inference.
The key distinctions are: \textsc{GLOBETROTTER} operates on a source--target
model rather than a full admixture graph topology, does not produce a graph
with edge probabilities, and uses LD-derived signals rather than IBD segments.
\hapgraph{} is therefore complementary: it provides a full graph topology and
IBD-based timing, while \textsc{GLOBETROTTER} provides richer source
decomposition and access to deeper time scales.

\subsection*{IBD-Based Demographic Inference}

The relationship between IBD segment lengths and demographic history was
formalised by \citet{Palamara2012}, who showed that the distribution of IBD
segment lengths encodes the coalescence rate function, enabling inference of
effective population size histories.
\citet{Browning2020} developed \hapibd{}, a scalable algorithm for phased IBD
detection that forms the pre-processing step in \hapgraph{}'s Stage II.
The use of mean IBD segment length as a clock for admixture timing, via the
$\mathbb{E}[\bar{L}] = 100/T$ model, is well-established in the IBD literature
\citep{Browning2020} and has been applied to date admixture events in diverse
human populations \citep{Moorjani2013}.
\hapgraph{} embeds this timing model within a full Bayesian inference
framework and integrates it with the F-statistic graph model, rather than
applying it as a standalone post-hoc analysis.

\subsection*{1000 Genomes Project and Human Population History}

\hapgraph{} was validated on the 2022 high-coverage 1kGP release
\citep{ByrskaBishop2022}, which provides phased SNP calls for 3202 individuals
across 26 populations and represents the highest-coverage large-scale human
population genomics reference to date.
The admixture events detected by \hapgraph{} are consistent with published
reconstructions of African American and Afro-Caribbean history
\citep{Bryc2015}, Latino population structure \citep{Moreno-Estrada2013},
and South Asian admixture chronology \citep{Moorjani2013}, providing
independent biological validation for the inferred graph.
